% GENERATED FILE. Edit canonical lessons, then run npm run generate:books.
% canonical-source-hash: fnv1a64:654de6c0bee90d42
% canonical-lessons: ES-C776-contrasena, ES-C776-archivo, ES-C776-carpeta, ES-C776-actualizar, ES-C776-teclado

\chapter{Si olvidas la contraseña}
\label{ch:es-si-olvidas-la-contrasena}
\hlchaptermodality{776}

\begin{chapteropening}
\textbf{By the end of this chapter:} \emph{I can say what to do if something goes wrong on a computer: a password, a file, a folder, an update or a keyboard.}

Complete the last lesson of chapter 776: I can say what to do if something goes wrong on a computer: a password, a file, a folder, an update or a keyboard.
\end{chapteropening}

\section[la contraseña]{la contraseña — the password}

\label{lesson:ES-C776-contrasena}



\begin{quote}
Before the new one: say the Spanish for the piece, then the Spanish for the smell.
\end{quote}

\subsection*{You'll want to know: la contraseña}
\textbf{la contraseña} — \textquotedblleft{}the password\textquotedblright{}.

\begin{quote}
Si olvidas la contraseña, pulsa aquí.
\end{quote}

\begin{quote}
\emph{If you forget your password, click here.}
\end{quote}

\emph{La clave} is also used, especially for a PIN.

\begin{cousinweb}
\emph{contraseña} is built on \emph{contra} 'against, in return' and \emph{seña} 'a sign', which comes from Latin \emph{signa}, the plural of \emph{signum} 'a mark' (Lewis \& Short). It was first a soldier's countersign.
\end{cousinweb}

\begin{grammarlens}[title={What you've built}]
One more word for telling, explaining and describing at B1.
\end{grammarlens}

\subsection*{Your turn}
\begin{itemize}
\raggedright
  \item \emph{Say it:} \emph{la contraseña}
  \item \emph{Say it:} \emph{la contraseña}, once more
  \item \emph{Say it:} a sentence of your own with \emph{la contraseña}
  \item \emph{Recall:} say the Spanish for the piece, then the Spanish for the smell, then say \emph{la contraseña} again
\end{itemize}

\begin{tcolorbox}[breakable,colback=teal!4,colframe=teal!35!black,title={Before you move on}]
What does la contraseña mean? (\textquotedblleft{}The password\textquotedblright{}.) Say it once more.
\end{tcolorbox}
\section[el archivo]{el archivo — the file; the archive}

\label{lesson:ES-C776-archivo}



\begin{quote}
Before the new one: say the Spanish for the smell, then the Spanish for the password.
\end{quote}

\subsection*{You'll want to know: el archivo}
\textbf{el archivo} — \textquotedblleft{}the file; the archive\textquotedblright{}.

\begin{quote}
Si el archivo es muy grande, no se puede enviar.
\end{quote}

\begin{quote}
\emph{If the file is very big, it can't be sent.}
\end{quote}

\emph{Un archivo adjunto} — an attachment.

\begin{cousinweb}
\emph{archivo} is a learned borrowing from Late Latin \emph{archīvum} 'the archives' (Lewis \& Short), which comes from Greek \emph{archeîon} 'a town hall' (EtymDB).
\end{cousinweb}

\begin{grammarlens}[title={What you've built}]
One more word for telling, explaining and describing at B1.
\end{grammarlens}

\subsection*{Your turn}
\begin{itemize}
\raggedright
  \item \emph{Say it:} \emph{el archivo}
  \item \emph{Say it:} \emph{el archivo}, once more
  \item \emph{Say it:} a sentence of your own with \emph{el archivo}
  \item \emph{Recall:} say the Spanish for the smell, then the Spanish for the password, then say \emph{el archivo} again
\end{itemize}

\begin{tcolorbox}[breakable,colback=teal!4,colframe=teal!35!black,title={Before you move on}]
What does el archivo mean? (\textquotedblleft{}The file; the archive\textquotedblright{}.) Say it once more.
\end{tcolorbox}
\section[la carpeta]{la carpeta — the folder}

\label{lesson:ES-C776-carpeta}



\begin{quote}
Before the new one: say the Spanish for the password, then the Spanish for the file.
\end{quote}

\subsection*{You'll want to know: la carpeta}
\textbf{la carpeta} — \textquotedblleft{}the folder\textquotedblright{}.

\begin{quote}
Si no encuentras la foto, mira en esta carpeta.
\end{quote}

\begin{quote}
\emph{If you can't find the photo, look in this folder.}
\end{quote}

Also a cardboard folder for papers.

\begin{cousinweb}
\emph{carpeta} was borrowed from French \emph{carpette} (EtymDB), a cloth cover for a table, which goes back to Medieval Latin \emph{carpīta}. The cover for papers became a folder.
\end{cousinweb}

\begin{grammarlens}[title={What you've built}]
One more word for telling, explaining and describing at B1.
\end{grammarlens}

\subsection*{Your turn}
\begin{itemize}
\raggedright
  \item \emph{Say it:} \emph{la carpeta}
  \item \emph{Say it:} \emph{la carpeta}, once more
  \item \emph{Say it:} a sentence of your own with \emph{la carpeta}
  \item \emph{Recall:} say the Spanish for the password, then the Spanish for the file, then say \emph{la carpeta} again
\end{itemize}

\begin{tcolorbox}[breakable,colback=teal!4,colframe=teal!35!black,title={Before you move on}]
What does la carpeta mean? (\textquotedblleft{}The folder\textquotedblright{}.) Say it once more.
\end{tcolorbox}
\section[actualizar]{actualizar — to update}

\label{lesson:ES-C776-actualizar}



\begin{quote}
Before the new one: say the Spanish for the file, then the Spanish for the folder.
\end{quote}

\subsection*{You'll want to know: actualizar}
\textbf{actualizar} — \textquotedblleft{}to update\textquotedblright{}.

\begin{quote}
Si no funciona, actualiza la aplicación.
\end{quote}

\begin{quote}
\emph{If it doesn't work, update the app.}
\end{quote}

It takes \emph{c} before \emph{e}: \emph{actualicé}. \emph{La actualización} is the update.

\begin{cousinweb}
\emph{actualizar} is built on \emph{actual} 'current' with the ending \emph{-izar}. \emph{Actual} comes from Latin \emph{āctuālis} 'active, practical' (EtymDB; Lewis \& Short), from \emph{agere} 'to do'. In Spanish \emph{actual} means current, not real.
\end{cousinweb}

\begin{grammarlens}[title={What you've built}]
One more word for telling, explaining and describing at B1.
\end{grammarlens}

\subsection*{Your turn}
\begin{itemize}
\raggedright
  \item \emph{Say it:} \emph{actualizar}
  \item \emph{Say it:} \emph{actualizar}, once more
  \item \emph{Say it:} a sentence of your own with \emph{actualizar}
  \item \emph{Recall:} say the Spanish for the file, then the Spanish for the folder, then say \emph{actualizar} again
\end{itemize}

\begin{tcolorbox}[breakable,colback=teal!4,colframe=teal!35!black,title={Before you move on}]
What does actualizar mean? (\textquotedblleft{}To update\textquotedblright{}.) Say it once more.
\end{tcolorbox}
\section[el teclado]{el teclado — the keyboard}

\label{lesson:ES-C776-teclado}



\begin{quote}
Before the new one: say the Spanish for the folder, then the Spanish for to update.
\end{quote}

\subsection*{You'll want to know: el teclado}
\textbf{el teclado} — \textquotedblleft{}the keyboard\textquotedblright{}.

\begin{quote}
Si el teclado no funciona, prueba con otro.
\end{quote}

\begin{quote}
\emph{If the keyboard doesn't work, try another one.}
\end{quote}

\emph{Una tecla} is a key.

\begin{cousinweb}
\emph{teclado} is built on \emph{tecla} 'a key' with the ending \emph{-ado} (EtymDB). The origin of \emph{tecla} itself is not settled.
\end{cousinweb}

\begin{grammarlens}[title={What you've built}]
One more word for telling, explaining and describing at B1.
\end{grammarlens}

\subsection*{Your turn}
\begin{itemize}
\raggedright
  \item \emph{Say it:} \emph{el teclado}
  \item \emph{Say it:} \emph{el teclado}, once more
  \item \emph{Say it:} a sentence of your own with \emph{el teclado}
  \item \emph{Recall:} say the Spanish for the folder, then the Spanish for to update, then say \emph{el teclado} again
\end{itemize}

\begin{tcolorbox}[breakable,colback=teal!4,colframe=teal!35!black,title={Before you move on}]
What does el teclado mean? (\textquotedblleft{}The keyboard\textquotedblright{}.) Say it once more.
\end{tcolorbox}
